\section{Introduction}
\label{sec:introduction}

Large language models (LLMs) have demonstrated strong capabilities in general-purpose code generation \cite{vaswani2017attention, openai2023gpt4, yang2025qwen3}. However, applying LLMs to EDA scripting—where engineers routinely manipulate in-memory design databases, create layout instances, route schematics, and query netlists—remains challenging. First, EDA APIs are long-tailed and tool-specific; many rarely appear in public training corpora. Second, correctness cannot be verified statically: it depends on observable side effects in a real EDA environment, such as modified layouts or updated schematics. Third, without execution feedback, static generation cannot diagnose errors or iteratively refine incorrect code \cite{chen2024dawning, blocklove2024eda}. Existing LLM-based EDA work has mainly focused on HDL generation \cite{thakur2023verigen, liu2023rtlcoder}, while scripting with real tool APIs—especially with incomplete documentation and sandbox execution—remains underexplored.
% 大语言模型在通用代码生成方面表现强大，但应用于EDA脚本——工程师日常操作内存中的设计数据库、创建版图实例、布线原理图、查询网表——仍面临挑战。第一，EDA API具有长尾性和工具特定性，许多很少出现在公开训练语料中。第二，正确性无法静态验证：它取决于真实EDA环境中的可观察副作用，如版图修改或原理图更新。第三，没有执行反馈，静态生成无法诊断错误或迭代修正代码。现有LLM-for-EDA工作主要集中在HDL生成，而面向真实工具API、沙盒执行和不完整文档的脚本生成仍研究不足。

This paper presents \textbf{ZhuLong}, an execution-grounded LLM coding agent for PyAether, SKILL, and Tcl. Named after a mythical creature whose opened eyes bring light, ZhuLong aims to illuminate opaque EDA APIs and automate tool-specific scripting workflows. Built on Cline-CLI \cite{cline2024cline}, ZhuLong extends three MCP tools \cite{lee2025autoeda}—\texttt{search\_apis}, \texttt{get\_api\_details}, and \texttt{run\_code}—that form a closed loop: retrieve candidate APIs, inspect their documentation, execute code in the EDA environment with execution feedback, and iteratively refine based on observed errors or outputs. All tools accept a unified \texttt{lang} parameter, so the same agent logic covers PyAether, SKILL, and Tcl. To further address incomplete or outdated documentation, ZhuLong introduces an offline \textbf{API self-exploration mechanism} that proactively explores undocumented API behaviors in the sandbox via \textbf{counterfactual experimentation}---deliberately testing alternative parameter values, observing execution outcomes, inferring constraints and usage patterns, and storing the discovered information as enhanced API documentation for future queries. Beyond the engineering, we provide a \textbf{theoretical foundation} (Section~\ref{sec:grounding-theory}): EDA scripting correctness is a \emph{situated} property that cannot be guaranteed from static knowledge, and a Bayesian analysis shows the resulting ``grounding gap'' is bounded by observation fidelity---not model scale---explaining why runtime observation is essential, and why the Harness must supply both halves of it: reading the environment's documented API facts, and reading back what executing code does to the live design.
% 本文提出\textbf{ZhuLong}，一个面向PyAether和SKILL的执行驱动LLM代码智能体。其名源自中国神话中睁眼带来光明的神兽，寓意照亮不透明的EDA API并自动化工具特定的脚本流程。ZhuLong基于Cline-CLI构建，扩展三个MCP工具——\texttt{search\_apis}、\texttt{get\_api\_details}和\texttt{run\_code}——形成闭环：检索候选API、查阅文档、在沙盒EDA环境中执行代码、基于观察到的错误或输出迭代修正。所有工具接受统一的\texttt{lang}参数，以相同智能体逻辑同时支持PyAether和SKILL。为了进一步解决文档不完整或过时的问题，ZhuLong引入了离线\textbf{API自探索机制}，通过沙盒中的\textbf{反事实实验}主动探索未文档化的API行为——有意识地测试不同参数取值、观察执行结果、推断约束和使用模式、将发现的信息存储为增强的API文档供后续查询使用。

The main contributions are:

\begin{itemize}
    \item \textbf{An account of \emph{which} factor sets the ceiling, with measurements that distinguish them (RQ1).}
    We formalize the \emph{grounding gap}: EDA correctness is a situated (L4) property that no
    offline knowledge can guarantee, and a Bayesian argument shows the resulting gap is
    bounded by \emph{observation fidelity} rather than by the capacity of the model
    (Section~\ref{sec:grounding-theory}). We answer RQ1 from both sides---laddering the
    observation channel with the LLM fixed, and changing the prior with the harness
    fixed---and the framework is stated \emph{before} the measurements, so that its claims are
    predictions the experiments could refute rather than explanations fitted afterwards.

    \item \textbf{A harness whose necessary components we isolate (RQ2).}
    ZhuLong couples three unified MCP tools---\texttt{search\_apis}, \texttt{get\_api\_details},
    and \texttt{run\_code}---so that API retrieval and documentation inspection (the
    \emph{observation channel}), sandbox execution (the \emph{effector}), and the
    generate--execute--refine loop (the \emph{belief update}) form one closed loop, driven by a
    single \texttt{lang} parameter so the same agent logic covers PyAether, SKILL, and Tcl
    (Section~\ref{sec:system}). In a commercial Empyrean Aether environment, ZhuLong attains
    \textbf{78.5\%} Pass@1 on the 158-task PyAether slice, versus \textbf{23.6\%} for a pure
    LLM and \textbf{32.3\%} for a retrieval-augmented static baseline, and \textbf{[X\%]}
    overall on the unified benchmark. % TODO: fill unified-benchmark Pass@1 once SKILL/Tcl slices are done
    Ablation shows that the three components are \emph{jointly necessary and not
    interchangeable}: adding sandbox execution to RAG yields \textbf{+43.0~pp}, while removing
    execution ($-41.2$~pp) or retrieval ($-44.3$~pp) collapses accuracy to $\sim$35\%, and
    bounding the belief update degrades it monotonically (Section~\ref{sec:experiments}). The
    offline \textbf{API self-exploration} mechanism, measured as one rung of the observation
    ladder, adds \textbf{+3.2~pp} and reduces tool calls by \textbf{22.1\%} per trace
    (Section~\ref{sec:selfdoc}); as prior-side enrichment it is bounded by observation rather
    than a substitute for it.

    \item \textbf{A cross-language benchmark on which the same structure can be re-measured per toolchain (RQ3).}
    \textbf{EDA-Scripting-Bench} places PyAether, SKILL and Tcl under one task schema, one
    assertion criterion and one harness, so that the structure of RQ1--RQ2 can be re-measured
    slice by slice with nothing changed but the toolchain (Section~\ref{sec:benchmark}). Its
    completed PyAether slice, \textbf{EDA-Eval-PyAether} (158 tasks from four sources,
    expert-reviewed), serves as the first standardized PyAether benchmark; the SKILL and
    Tcl slices complete the suite for reproducible comparison of LLM-based scripting agents.
    What we release is the \textbf{harness, the task schema and the evaluation protocol}; the
    task cases themselves derive from vendor documentation and customer designs and cannot be
    redistributed.
\end{itemize}